\chapter{Core Code Listings}

\section{Purpose}
This appendix preserves representative implementation excerpts from the SARS project in a LaTeX-safe form. The goal is to document the most important engineering ideas without depending on external source-file imports during compilation.

\section{Database Initialization Excerpt}
\begin{lstlisting}[style=sarscode,caption={Representative backend database setup excerpt}]
from sqlalchemy import create_engine
from sqlalchemy.orm import declarative_base, sessionmaker

engine = create_engine(settings.DATABASE_URL, connect_args={"check_same_thread": False})
SessionLocal = sessionmaker(autocommit=False, autoflush=False, bind=engine)
Base = declarative_base()

def get_db():
    db = SessionLocal()
    try:
        yield db
    finally:
        db.close()
\end{lstlisting}

\section{Representative ORM Model Excerpt}
\begin{lstlisting}[style=sarscode,caption={Representative academic ORM model excerpt}]
class SemesterRecord(Base):
    __tablename__ = "semester_records"

    id = Column(Integer, primary_key=True, index=True)
    student_id = Column(Integer, ForeignKey("students.id"), nullable=False, index=True)
    semester_no = Column(Integer, nullable=False)
    gpa = Column(Float, nullable=True)
    credits_attempted = Column(Float, nullable=True)
    credits_earned = Column(Float, nullable=True)
    uploaded_at = Column(DateTime(timezone=True), server_default=func.now())
    is_confirmed = Column(Boolean, default=False)

class SubjectGrade(Base):
    __tablename__ = "subject_grades"

    id = Column(Integer, primary_key=True, index=True)
    semester_record_id = Column(Integer, ForeignKey("semester_records.id"), nullable=False)
    subject_name = Column(String, nullable=False)
    subject_code = Column(String, nullable=True)
    grade_letter = Column(String, nullable=True)
    grade_points = Column(Float, nullable=True)
    credits = Column(Float, nullable=True)
    is_backlog = Column(Boolean, default=False)
\end{lstlisting}

\section{Risk Engine Excerpt}
\begin{lstlisting}[style=sarscode,caption={Representative SARS risk computation excerpt}]
def compute_cgpa(student_id: int, db: Session):
    records = db.query(SemesterRecord).filter(
        SemesterRecord.student_id == student_id
    ).all()

    total_weighted_points = 0.0
    total_credits = 0.0

    for rec in records:
        if rec.gpa is None or rec.credits_attempted in (None, 0):
            continue
        total_weighted_points += float(rec.gpa) * float(rec.credits_attempted)
        total_credits += float(rec.credits_attempted)

    if total_credits == 0:
        return None

    return round(total_weighted_points / total_credits, 2)

def final_sars_score(gpa_risk, backlog_risk, attendance_risk):
    return round(
        gpa_risk * 0.40 +
        backlog_risk * 0.35 +
        attendance_risk * 0.25,
        2
    )
\end{lstlisting}

\section{RAG Chunk Construction Excerpt}
\begin{lstlisting}[style=sarscode,caption={Representative student-specific RAG chunk construction excerpt}]
def build_risk_summary_chunk(student, risk, breakdown):
    lines = [
        "Risk summary:",
        f"- SARS score: {risk.sars_score}",
        f"- Risk level: {risk.risk_level}",
        f"- Confidence: {risk.confidence}",
    ]

    for key in sorted(breakdown.keys()):
        lines.append(f"- {key}: {breakdown[key]}")

    return {
        "chunk_type": "risk_summary",
        "student_id": str(student.id),
        "text": "\\n".join(lines),
    }
\end{lstlisting}

\section{Interpretation}
These excerpts are sufficient to document the key engineering contribution of the system:
\begin{itemize}[leftmargin=*]
\item the relational backend is session-based and modular,
\item academic records are normalized into semester and subject entities,
\item risk logic is transparent and formula-driven,
\item retrieval content is built from student-specific academic evidence.
\end{itemize}
